\documentclass[letterpaper,12pt]{article}
\usepackage[margin=.5in]{geometry}
\usepackage{tikz}
\usepackage{helvet}
\renewcommand{\familydefault}{\sfdefault}
\pagestyle{empty}
\pdfmapfile{+cm.map}
\pdfmapfile{+ps2pk35.map}
\setlength{\parindent}{0pt}
\hyphenpenalty=10000
\exhyphenpenalty=10000
\begin{document}
\null\begin{tikzpicture}[remember picture,overlay,x=1cm,y=1cm]
\begin{scope}[shift={(current page.north west)}]
\node[anchor=north west,inner sep=0pt,font=\fontsize{11.5}{13}\selectfont] at (1.4,-.72) {Week 29 / Two-length builders / Grades 4--5};
\node[anchor=south west,inner sep=0pt,font=\fontsize{9}{11}\selectfont] at (1.4,-27.2) {Bellingham Math Circle / Week 29 / F29-45-v2};
\node[anchor=south east,inner sep=0pt,font=\fontsize{9}{11}\selectfont] at (20.15,-27.2) {1};
\begin{scope}[shift={(1.4,-1.65)}]
\node[anchor=north west,inner sep=0pt,align=left,text width=18.5cm,font=\fontsize{13}{16.38}\selectfont] at (0,0) {Put whole rods end to end. Leave no gaps or overlaps. You may use a length more than once. Changing the order of the same rods does not make a new way.};
\draw[fill=gray!12,line width=.7pt] (0,-2) rectangle (3,-3);
\draw[gray!60,line width=.3pt] (1,-2) -- (1,-3);
\draw[gray!60,line width=.3pt] (2,-2) -- (2,-3);
\node[font=\fontsize{14}{16}\selectfont,inner sep=1pt] at (1.5,-2.5) {3};
\draw[fill=gray!12,line width=.7pt] (3.9,-2) rectangle (7.9,-3);
\draw[gray!60,line width=.3pt] (4.9,-2) -- (4.9,-3);
\draw[gray!60,line width=.3pt] (5.9,-2) -- (5.9,-3);
\draw[gray!60,line width=.3pt] (6.9,-2) -- (6.9,-3);
\node[font=\fontsize{14}{16}\selectfont,inner sep=1pt] at (5.9,-2.5) {4};
\draw[line width=.7pt] (0,-3.6) rectangle (7,-4.6);
\draw[gray!60,line width=.3pt] (1,-3.6) -- (1,-4.6);
\draw[gray!60,line width=.3pt] (2,-3.6) -- (2,-4.6);
\draw[gray!60,line width=.3pt] (3,-3.6) -- (3,-4.6);
\draw[gray!60,line width=.3pt] (4,-3.6) -- (4,-4.6);
\draw[gray!60,line width=.3pt] (5,-3.6) -- (5,-4.6);
\draw[gray!60,line width=.3pt] (6,-3.6) -- (6,-4.6);
\draw[line width=1.5pt] (0,-3.6) rectangle (3,-4.6);
\draw[line width=1.5pt] (3,-3.6) rectangle (7,-4.6);
\node[font=\fontsize{12}{14}\selectfont,inner sep=1pt] at (1.5,-4.1) {3};
\node[font=\fontsize{12}{14}\selectfont,inner sep=1pt] at (5,-4.1) {4};
\node[font=\fontsize{14}{16}\selectfont,inner sep=1pt] at (9,-4.1) {3 + 4 = 7};
\node[anchor=north west,inner sep=0pt,align=left,text width=18.6cm,font=\fontsize{14}{17.78}\selectfont] at (0,-5.7) {\textbf{Problem 1:} For 3-rods and 4-rods, find every impossible positive length. Explain why your list is complete.};
\draw[black,line width=.65pt] (0.0,-8.0) rectangle (1.8,-9.2);
\node[font=\fontsize{16}{18}\selectfont,inner sep=1pt] at (0.9,-8.6) {1};
\draw[black,line width=.65pt] (2.2,-8.0) rectangle (4.0,-9.2);
\node[font=\fontsize{16}{18}\selectfont,inner sep=1pt] at (3.1,-8.6) {2};
\draw[black,line width=.65pt] (4.4,-8.0) rectangle (6.2,-9.2);
\node[font=\fontsize{16}{18}\selectfont,inner sep=1pt] at (5.300000000000001,-8.6) {3};
\draw[black,line width=.65pt] (6.6000000000000005,-8.0) rectangle (8.4,-9.2);
\node[font=\fontsize{16}{18}\selectfont,inner sep=1pt] at (7.500000000000001,-8.6) {4};
\draw[black,line width=.65pt] (8.8,-8.0) rectangle (10.600000000000001,-9.2);
\node[font=\fontsize{16}{18}\selectfont,inner sep=1pt] at (9.700000000000001,-8.6) {5};
\draw[black,line width=.65pt] (11.0,-8.0) rectangle (12.8,-9.2);
\node[font=\fontsize{16}{18}\selectfont,inner sep=1pt] at (11.9,-8.6) {6};
\draw[black,line width=.65pt] (13.200000000000001,-8.0) rectangle (15.000000000000002,-9.2);
\node[font=\fontsize{16}{18}\selectfont,inner sep=1pt] at (14.100000000000001,-8.6) {7};
\draw[black,line width=.65pt] (0.0,-9.6) rectangle (1.8,-10.799999999999999);
\node[font=\fontsize{16}{18}\selectfont,inner sep=1pt] at (0.9,-10.2) {8};
\draw[black,line width=.65pt] (2.2,-9.6) rectangle (4.0,-10.799999999999999);
\node[font=\fontsize{16}{18}\selectfont,inner sep=1pt] at (3.1,-10.2) {9};
\draw[black,line width=.65pt] (4.4,-9.6) rectangle (6.2,-10.799999999999999);
\node[font=\fontsize{16}{18}\selectfont,inner sep=1pt] at (5.300000000000001,-10.2) {10};
\draw[black,line width=.65pt] (6.6000000000000005,-9.6) rectangle (8.4,-10.799999999999999);
\node[font=\fontsize{16}{18}\selectfont,inner sep=1pt] at (7.500000000000001,-10.2) {11};
\draw[black,line width=.65pt] (8.8,-9.6) rectangle (10.600000000000001,-10.799999999999999);
\node[font=\fontsize{16}{18}\selectfont,inner sep=1pt] at (9.700000000000001,-10.2) {12};
\draw[black,line width=.65pt] (11.0,-9.6) rectangle (12.8,-10.799999999999999);
\node[font=\fontsize{16}{18}\selectfont,inner sep=1pt] at (11.9,-10.2) {13};
\draw[black,line width=.65pt] (13.200000000000001,-9.6) rectangle (15.000000000000002,-10.799999999999999);
\node[font=\fontsize{16}{18}\selectfont,inner sep=1pt] at (14.100000000000001,-10.2) {14};
\draw[black,line width=.65pt] (0.0,-11.2) rectangle (1.8,-12.399999999999999);
\node[font=\fontsize{16}{18}\selectfont,inner sep=1pt] at (0.9,-11.799999999999999) {15};
\draw[black,line width=.65pt] (2.2,-11.2) rectangle (4.0,-12.399999999999999);
\node[font=\fontsize{16}{18}\selectfont,inner sep=1pt] at (3.1,-11.799999999999999) {16};
\draw[black,line width=.65pt] (4.4,-11.2) rectangle (6.2,-12.399999999999999);
\node[font=\fontsize{16}{18}\selectfont,inner sep=1pt] at (5.300000000000001,-11.799999999999999) {17};
\draw[black,line width=.65pt] (6.6000000000000005,-11.2) rectangle (8.4,-12.399999999999999);
\node[font=\fontsize{16}{18}\selectfont,inner sep=1pt] at (7.500000000000001,-11.799999999999999) {18};
\draw[black,line width=.65pt] (8.8,-11.2) rectangle (10.600000000000001,-12.399999999999999);
\node[font=\fontsize{16}{18}\selectfont,inner sep=1pt] at (9.700000000000001,-11.799999999999999) {19};
\draw[black,line width=.65pt] (11.0,-11.2) rectangle (12.8,-12.399999999999999);
\node[font=\fontsize{16}{18}\selectfont,inner sep=1pt] at (11.9,-11.799999999999999) {20};
\draw[gray!45,line width=.35pt] (0,-14.0) -- (18.5,-14.0);
\draw[gray!45,line width=.35pt] (0,-15.15) -- (18.5,-15.15);
\draw[gray!45,line width=.35pt] (0,-16.3) -- (18.5,-16.3);
\draw[gray!45,line width=.35pt] (0,-17.45) -- (18.5,-17.45);
\draw[gray!45,line width=.35pt] (0,-18.6) -- (18.5,-18.6);
\draw[gray!45,line width=.35pt] (0,-19.75) -- (18.5,-19.75);
\draw[gray!45,line width=.35pt] (0,-20.9) -- (18.5,-20.9);
\draw[gray!45,line width=.35pt] (0,-22.049999999999997) -- (18.5,-22.049999999999997);
\draw[gray!45,line width=.35pt] (0,-23.2) -- (18.5,-23.2);\end{scope}\end{scope}\end{tikzpicture}
\newpage
\null\begin{tikzpicture}[remember picture,overlay,x=1cm,y=1cm]
\begin{scope}[shift={(current page.north west)}]
\node[anchor=north west,inner sep=0pt,font=\fontsize{11.5}{13}\selectfont] at (1.4,-.72) {Week 29 / Two-length builders / Grades 4--5};
\node[anchor=south west,inner sep=0pt,font=\fontsize{9}{11}\selectfont] at (1.4,-27.2) {Bellingham Math Circle / Week 29 / F29-45-v2};
\node[anchor=south east,inner sep=0pt,font=\fontsize{9}{11}\selectfont] at (20.15,-27.2) {2};
\begin{scope}[shift={(1.4,-1.65)}]
\node[anchor=north west,inner sep=0pt,align=left,text width=18.6cm,font=\fontsize{14}{17.78}\selectfont] at (0,0) {\textbf{Problem 2:} Find the last impossible length for each pair: 3 and 5; 4 and 5; 4 and 7. For each pair, give a finite collection of builds that settles every longer target.};
\node[anchor=north west,inner sep=0pt,align=left,text width=18.5cm,font=\fontsize{14}{17.64}\selectfont] at (0,-2.4) {3 and 5};
\draw[gray!45,line width=.35pt] (0,-4.0) -- (18.5,-4.0);
\draw[gray!45,line width=.35pt] (0,-5.15) -- (18.5,-5.15);
\draw[gray!45,line width=.35pt] (0,-6.3) -- (18.5,-6.3);
\draw[gray!45,line width=.35pt] (0,-7.449999999999999) -- (18.5,-7.449999999999999);
\node[anchor=north west,inner sep=0pt,align=left,text width=18.5cm,font=\fontsize{14}{17.64}\selectfont] at (0,-9.4) {4 and 5};
\draw[gray!45,line width=.35pt] (0,-11.0) -- (18.5,-11.0);
\draw[gray!45,line width=.35pt] (0,-12.15) -- (18.5,-12.15);
\draw[gray!45,line width=.35pt] (0,-13.3) -- (18.5,-13.3);
\draw[gray!45,line width=.35pt] (0,-14.45) -- (18.5,-14.45);
\node[anchor=north west,inner sep=0pt,align=left,text width=18.5cm,font=\fontsize{14}{17.64}\selectfont] at (0,-16.4) {4 and 7};
\draw[gray!45,line width=.35pt] (0,-18.0) -- (18.5,-18.0);
\draw[gray!45,line width=.35pt] (0,-19.15) -- (18.5,-19.15);
\draw[gray!45,line width=.35pt] (0,-20.3) -- (18.5,-20.3);
\draw[gray!45,line width=.35pt] (0,-21.45) -- (18.5,-21.45);\end{scope}\end{scope}\end{tikzpicture}
\newpage
\null\begin{tikzpicture}[remember picture,overlay,x=1cm,y=1cm]
\begin{scope}[shift={(current page.north west)}]
\node[anchor=north west,inner sep=0pt,font=\fontsize{11.5}{13}\selectfont] at (1.4,-.72) {Week 29 / Two-length builders / Grades 4--5};
\node[anchor=south west,inner sep=0pt,font=\fontsize{9}{11}\selectfont] at (1.4,-27.2) {Bellingham Math Circle / Week 29 / F29-45-v2};
\node[anchor=south east,inner sep=0pt,font=\fontsize{9}{11}\selectfont] at (20.15,-27.2) {3};
\begin{scope}[shift={(1.4,-1.65)}]
\node[anchor=north west,inner sep=0pt,align=left,text width=18.6cm,font=\fontsize{14}{17.78}\selectfont] at (0,0) {\textbf{Problem 3:} Which lengths from 18 through 29 can be built with 4-rods and 7-rods? Convince someone without checking twelve separate cases.};
\draw[black,line width=.65pt] (0.0,-3.0) rectangle (1.8,-4.2);
\node[font=\fontsize{16}{18}\selectfont,inner sep=1pt] at (0.9,-3.6) {18};
\draw[black,line width=.65pt] (2.2,-3.0) rectangle (4.0,-4.2);
\node[font=\fontsize{16}{18}\selectfont,inner sep=1pt] at (3.1,-3.6) {19};
\draw[black,line width=.65pt] (4.4,-3.0) rectangle (6.2,-4.2);
\node[font=\fontsize{16}{18}\selectfont,inner sep=1pt] at (5.300000000000001,-3.6) {20};
\draw[black,line width=.65pt] (6.6000000000000005,-3.0) rectangle (8.4,-4.2);
\node[font=\fontsize{16}{18}\selectfont,inner sep=1pt] at (7.500000000000001,-3.6) {21};
\draw[black,line width=.65pt] (8.8,-3.0) rectangle (10.600000000000001,-4.2);
\node[font=\fontsize{16}{18}\selectfont,inner sep=1pt] at (9.700000000000001,-3.6) {22};
\draw[black,line width=.65pt] (11.0,-3.0) rectangle (12.8,-4.2);
\node[font=\fontsize{16}{18}\selectfont,inner sep=1pt] at (11.9,-3.6) {23};
\draw[black,line width=.65pt] (0.0,-4.6) rectangle (1.8,-5.8);
\node[font=\fontsize{16}{18}\selectfont,inner sep=1pt] at (0.9,-5.199999999999999) {24};
\draw[black,line width=.65pt] (2.2,-4.6) rectangle (4.0,-5.8);
\node[font=\fontsize{16}{18}\selectfont,inner sep=1pt] at (3.1,-5.199999999999999) {25};
\draw[black,line width=.65pt] (4.4,-4.6) rectangle (6.2,-5.8);
\node[font=\fontsize{16}{18}\selectfont,inner sep=1pt] at (5.300000000000001,-5.199999999999999) {26};
\draw[black,line width=.65pt] (6.6000000000000005,-4.6) rectangle (8.4,-5.8);
\node[font=\fontsize{16}{18}\selectfont,inner sep=1pt] at (7.500000000000001,-5.199999999999999) {27};
\draw[black,line width=.65pt] (8.8,-4.6) rectangle (10.600000000000001,-5.8);
\node[font=\fontsize{16}{18}\selectfont,inner sep=1pt] at (9.700000000000001,-5.199999999999999) {28};
\draw[black,line width=.65pt] (11.0,-4.6) rectangle (12.8,-5.8);
\node[font=\fontsize{16}{18}\selectfont,inner sep=1pt] at (11.9,-5.199999999999999) {29};
\draw[gray!45,line width=.35pt] (0,-8.0) -- (18.5,-8.0);
\draw[gray!45,line width=.35pt] (0,-9.15) -- (18.5,-9.15);
\draw[gray!45,line width=.35pt] (0,-10.3) -- (18.5,-10.3);
\draw[gray!45,line width=.35pt] (0,-11.45) -- (18.5,-11.45);
\draw[gray!45,line width=.35pt] (0,-12.6) -- (18.5,-12.6);
\draw[gray!45,line width=.35pt] (0,-13.75) -- (18.5,-13.75);
\node[anchor=north west,inner sep=0pt,align=left,text width=18.6cm,font=\fontsize{14}{17.78}\selectfont] at (0,-17) {\textbf{Problem 4:} Can 4-rods and 7-rods make 17? Explain your answer.};
\draw[gray!45,line width=.35pt] (0,-20.0) -- (18.5,-20.0);
\draw[gray!45,line width=.35pt] (0,-21.15) -- (18.5,-21.15);
\draw[gray!45,line width=.35pt] (0,-22.3) -- (18.5,-22.3);\end{scope}\end{scope}\end{tikzpicture}
\newpage
\null\begin{tikzpicture}[remember picture,overlay,x=1cm,y=1cm]
\begin{scope}[shift={(current page.north west)}]
\node[anchor=north west,inner sep=0pt,font=\fontsize{11.5}{13}\selectfont] at (1.4,-.72) {Week 29 / Two-length builders / Grades 4--5};
\node[anchor=south west,inner sep=0pt,font=\fontsize{9}{11}\selectfont] at (1.4,-27.2) {Bellingham Math Circle / Week 29 / F29-45-v2};
\node[anchor=south east,inner sep=0pt,font=\fontsize{9}{11}\selectfont] at (20.15,-27.2) {4};
\begin{scope}[shift={(1.4,-1.65)}]
\node[anchor=north west,inner sep=0pt,align=left,text width=18.6cm,font=\fontsize{14}{17.78}\selectfont] at (0,0) {\textbf{Problem 5:} Which pairs have infinitely many impossible lengths: 2 and 4; 3 and 6; 4 and 6; 3 and 5; 4 and 7? Find a rule that separates them and explain it.};
\draw[gray!45,line width=.35pt] (0,-3.0) -- (18.5,-3.0);
\draw[gray!45,line width=.35pt] (0,-4.15) -- (18.5,-4.15);
\draw[gray!45,line width=.35pt] (0,-5.3) -- (18.5,-5.3);
\draw[gray!45,line width=.35pt] (0,-6.449999999999999) -- (18.5,-6.449999999999999);
\draw[gray!45,line width=.35pt] (0,-7.6) -- (18.5,-7.6);
\draw[gray!45,line width=.35pt] (0,-8.75) -- (18.5,-8.75);
\draw[gray!45,line width=.35pt] (0,-9.899999999999999) -- (18.5,-9.899999999999999);
\draw[gray!45,line width=.35pt] (0,-11.049999999999999) -- (18.5,-11.049999999999999);
\node[anchor=north west,inner sep=0pt,align=left,text width=18.6cm,font=\fontsize{14}{17.78}\selectfont] at (0,-14) {\textbf{Problem 6:} Find two different ways to build 24 with 3-rods and 4-rods, then find all the ways. Which replacement of rods changes one way into another?};
\draw[gray!45,line width=.35pt] (0,-17.0) -- (18.5,-17.0);
\draw[gray!45,line width=.35pt] (0,-18.15) -- (18.5,-18.15);
\draw[gray!45,line width=.35pt] (0,-19.3) -- (18.5,-19.3);
\draw[gray!45,line width=.35pt] (0,-20.45) -- (18.5,-20.45);
\draw[gray!45,line width=.35pt] (0,-21.6) -- (18.5,-21.6);
\draw[gray!45,line width=.35pt] (0,-22.75) -- (18.5,-22.75);\end{scope}\end{scope}\end{tikzpicture}
\newpage
\null\begin{tikzpicture}[remember picture,overlay,x=1cm,y=1cm]
\begin{scope}[shift={(current page.north west)}]
\node[anchor=north west,inner sep=0pt,font=\fontsize{11.5}{13}\selectfont] at (1.4,-.72) {Week 29 / Two-length builders / Grades 4--5};
\node[anchor=south west,inner sep=0pt,font=\fontsize{9}{11}\selectfont] at (1.4,-27.2) {Bellingham Math Circle / Week 29 / F29-45-v2};
\node[anchor=south east,inner sep=0pt,font=\fontsize{9}{11}\selectfont] at (20.15,-27.2) {5};
\begin{scope}[shift={(1.4,-1.65)}]
\node[anchor=north west,inner sep=0pt,align=left,text width=18.6cm,font=\fontsize{14}{17.78}\selectfont] at (0,0) {\textbf{Problem 7:} Predict the last impossible length for 5-rods and 7-rods, then test your prediction. Compare the answer with your earlier pairs and propose a rule for pairs that eventually have no gaps.};
\draw[black,line width=.65pt] (0.0,-3.0) rectangle (1.8,-4.2);
\node[font=\fontsize{16}{18}\selectfont,inner sep=1pt] at (0.9,-3.6) {19};
\draw[black,line width=.65pt] (2.2,-3.0) rectangle (4.0,-4.2);
\node[font=\fontsize{16}{18}\selectfont,inner sep=1pt] at (3.1,-3.6) {20};
\draw[black,line width=.65pt] (4.4,-3.0) rectangle (6.2,-4.2);
\node[font=\fontsize{16}{18}\selectfont,inner sep=1pt] at (5.300000000000001,-3.6) {21};
\draw[black,line width=.65pt] (6.6000000000000005,-3.0) rectangle (8.4,-4.2);
\node[font=\fontsize{16}{18}\selectfont,inner sep=1pt] at (7.500000000000001,-3.6) {22};
\draw[black,line width=.65pt] (8.8,-3.0) rectangle (10.600000000000001,-4.2);
\node[font=\fontsize{16}{18}\selectfont,inner sep=1pt] at (9.700000000000001,-3.6) {23};
\draw[black,line width=.65pt] (11.0,-3.0) rectangle (12.8,-4.2);
\node[font=\fontsize{16}{18}\selectfont,inner sep=1pt] at (11.9,-3.6) {24};
\draw[black,line width=.65pt] (13.200000000000001,-3.0) rectangle (15.000000000000002,-4.2);
\node[font=\fontsize{16}{18}\selectfont,inner sep=1pt] at (14.100000000000001,-3.6) {25};
\draw[black,line width=.65pt] (0.0,-4.6) rectangle (1.8,-5.8);
\node[font=\fontsize{16}{18}\selectfont,inner sep=1pt] at (0.9,-5.199999999999999) {26};
\draw[black,line width=.65pt] (2.2,-4.6) rectangle (4.0,-5.8);
\node[font=\fontsize{16}{18}\selectfont,inner sep=1pt] at (3.1,-5.199999999999999) {27};
\draw[black,line width=.65pt] (4.4,-4.6) rectangle (6.2,-5.8);
\node[font=\fontsize{16}{18}\selectfont,inner sep=1pt] at (5.300000000000001,-5.199999999999999) {28};
\draw[black,line width=.65pt] (6.6000000000000005,-4.6) rectangle (8.4,-5.8);
\node[font=\fontsize{16}{18}\selectfont,inner sep=1pt] at (7.500000000000001,-5.199999999999999) {29};
\draw[black,line width=.65pt] (8.8,-4.6) rectangle (10.600000000000001,-5.8);
\node[font=\fontsize{16}{18}\selectfont,inner sep=1pt] at (9.700000000000001,-5.199999999999999) {30};
\draw[black,line width=.65pt] (11.0,-4.6) rectangle (12.8,-5.8);
\node[font=\fontsize{16}{18}\selectfont,inner sep=1pt] at (11.9,-5.199999999999999) {31};
\draw[gray!45,line width=.35pt] (0,-8.0) -- (18.5,-8.0);
\draw[gray!45,line width=.35pt] (0,-9.15) -- (18.5,-9.15);
\draw[gray!45,line width=.35pt] (0,-10.3) -- (18.5,-10.3);
\draw[gray!45,line width=.35pt] (0,-11.45) -- (18.5,-11.45);
\draw[gray!45,line width=.35pt] (0,-12.6) -- (18.5,-12.6);
\node[anchor=north west,inner sep=0pt,align=left,text width=18.6cm,font=\fontsize{14}{17.78}\selectfont] at (0,-16) {\textbf{Problem 8:} Invent paper rods of two different whole-number lengths a and b that have no common divisor larger than 1. Explain why some finite group of builds must settle every sufficiently long target.};
\draw[gray!45,line width=.35pt] (0,-19.0) -- (18.5,-19.0);
\draw[gray!45,line width=.35pt] (0,-20.15) -- (18.5,-20.15);
\draw[gray!45,line width=.35pt] (0,-21.3) -- (18.5,-21.3);
\draw[gray!45,line width=.35pt] (0,-22.45) -- (18.5,-22.45);\end{scope}\end{scope}\end{tikzpicture}
\end{document}
